% GENERATED FILE. Edit canonical lessons, then run npm run generate:books.
% canonical-source-hash: fnv1a64:f0770122688399f0
% canonical-lessons: KA-C445-park-madu, KA-C445-tirugu, KA-C445-alugadu, KA-C445-calisu, KA-C445-samipisu

\chapter{Parking and Turning}
\label{ch:ka-parking-and-turning}
\hlchaptermodality{445}

\begin{chapteropening}
\textbf{By the end of this chapter:} \emph{I can say to park a vehicle, to turn round, to wobble, to move and to approach.}

Complete the last lesson of chapter 445: I can say to park a vehicle, to turn round, to wobble, to move and to approach.
\end{chapteropening}

\section[\texorpdfstring{\kn{ಪಾರ್ಕ್} \kn{ಮಾಡು} (pārk māḍu)}{ಪಾರ್ಕ್ ಮಾಡು (pārk māḍu)}]{\kn{ಪಾರ್ಕ್} \kn{ಮಾಡು} (pārk māḍu) — to park a vehicle; to park}

\label{lesson:KA-C445-park-madu}



\begin{quote}
Before the new one: say the Kannada for to ban, then the Kannada for one-way.
\end{quote}

\subsection*{You'll want to know: \kn{ಪಾರ್ಕ್} \kn{ಮಾಡು}}
\textbf{\kn{ಪಾರ್ಕ್} \kn{ಮಾಡು}} — \emph{pārk māḍu} — \textquotedblleft{}to park a vehicle; to park\textquotedblright{}.

\begin{quote}
\kn{ನಾನು} \kn{ಕಾರನ್ನು} \kn{ಅಂಗಡಿಯ} \kn{ಎದುರು} \kn{ಪಾರ್ಕ್} \kn{ಮಾಡಿದೆ}.
\end{quote}

\begin{quote}
\emph{nānu kārannu aṅgaḍiya eduru pārk māḍide.}
\end{quote}

\begin{quote}
I parked the car opposite the shop.
\end{quote}

Past (I …): \kn{ಪಾರ್ಕ್} \kn{ಮಾಡಿದೆ}, \emph{pārk māḍide}.

The vehicle takes -\kn{ಅನ್ನು}; \kn{ಪಾರ್ಕಿಂಗ್} is a parking place.

\begin{cousinweb}
\textbf{\kn{ಪಾರ್ಕ್} \kn{ಮಾಡು}} (\emph{pārk māḍu}) joins a loan from English \emph{park} and the native verb \kn{ಮಾಡು}, to do.
\end{cousinweb}

\begin{grammarlens}[title={What you've built}]
One more word for getting around.
\end{grammarlens}

\subsection*{Your turn}
\begin{itemize}
\raggedright
  \item \emph{Say it:} \emph{pārk māḍu}
  \item \emph{Say it:} \emph{pārk māḍu}, once more
  \item \emph{Say it:} \emph{nānu kārannu aṅgaḍiya eduru pārk māḍide.} — I parked the car opposite the shop.
  \item \emph{Recall:} say the Kannada for to ban, then the Kannada for one-way, then say \emph{pārk māḍu} again
\end{itemize}

\begin{tcolorbox}[breakable,colback=teal!4,colframe=teal!35!black,title={Before you move on}]
What does \kn{ಪಾರ್ಕ್} \kn{ಮಾಡು} mean? (\textquotedblleft{}To park a vehicle; to park\textquotedblright{}.) Say it once more.
\end{tcolorbox}
\section[\texorpdfstring{\kn{ತಿರುಗು} (tirugu)}{ತಿರುಗು (tirugu)}]{\kn{ತಿರುಗು} (tirugu) — to turn round; to go round}

\label{lesson:KA-C445-tirugu}



\begin{quote}
Before the new one: say the Kannada for one-way, then the Kannada for to park a vehicle.
\end{quote}

\subsection*{You'll want to know: \kn{ತಿರುಗು}}
\textbf{\kn{ತಿರುಗು}} — \emph{tirugu} — \textquotedblleft{}to turn round; to go round\textquotedblright{}.

\begin{quote}
\kn{ಮುಂದಿನ} \kn{ವೃತ್ತದಲ್ಲಿ} \kn{ಬಸ್} \kn{ಎಡಕ್ಕೆ} \kn{ತಿರುಗುತ್ತದೆ}.
\end{quote}

\begin{quote}
\emph{mundina vṛttadalli bas eḍakke tiruguttade.}
\end{quote}

\begin{quote}
At the next circle the bus turns left.
\end{quote}

Past (I …): \kn{ತಿರುಗಿದೆ}, \emph{tirugide}.

\kn{ಹಿಂತಿರುಗು}, to come back, is \kn{ಹಿಂ}-, back, plus this verb.

\begin{cousinweb}
\textbf{\kn{ತಿರುಗು}} (\emph{tirugu}) is native Dravidian: DEDR 3246 lists Kannada \emph{tirigu}, \emph{tirugu}, \textquotedblleft{}to go round, to turn, to roll, to go back\textquotedblright{}.
\end{cousinweb}

\begin{grammarlens}[title={What you've built}]
One more word for getting around.
\end{grammarlens}

\subsection*{Your turn}
\begin{itemize}
\raggedright
  \item \emph{Say it:} \emph{tirugu}
  \item \emph{Say it:} \emph{tirugu}, once more
  \item \emph{Say it:} \emph{mundina vṛttadalli bas eḍakke tiruguttade.} — At the next circle the bus turns left.
  \item \emph{Recall:} say the Kannada for one-way, then the Kannada for to park a vehicle, then say \emph{tirugu} again
\end{itemize}

\begin{tcolorbox}[breakable,colback=teal!4,colframe=teal!35!black,title={Before you move on}]
What does \kn{ತಿರುಗು} mean? (\textquotedblleft{}To turn round; to go round\textquotedblright{}.) Say it once more.
\end{tcolorbox}
\section[\texorpdfstring{\kn{ಅಲುಗಾಡು} (alugāḍu)}{ಅಲುಗಾಡು (alugāḍu)}]{\kn{ಅಲುಗಾಡು} (alugāḍu) — to wobble; to shake}

\label{lesson:KA-C445-alugadu}



\begin{quote}
Before the new one: say the Kannada for to park a vehicle, then the Kannada for to turn round.
\end{quote}

\subsection*{You'll want to know: \kn{ಅಲುಗಾಡು}}
\textbf{\kn{ಅಲುಗಾಡು}} — \emph{alugāḍu} — \textquotedblleft{}to wobble; to shake\textquotedblright{}.

\begin{quote}
\kn{ಹಳೆಯ} \kn{ಬಸ್} \kn{ಕಲ್ಲಿನ} \kn{ರಸ್ತೆಯಲ್ಲಿ} \kn{ಅಲುಗಾಡುತ್ತಿತ್ತು}.
\end{quote}

\begin{quote}
\emph{haḷeya bas kallina rasteyalli alugāḍuttittu.}
\end{quote}

\begin{quote}
The old bus was rattling along the stony road.
\end{quote}

Past (it …): \kn{ಅಲುಗಾಡಿತು}, \emph{alugāḍitu}.

It shakes by itself; \kn{ಅಲ್ಲಾಡಿಸು} is to shake something.

\begin{cousinweb}
\textbf{\kn{ಅಲುಗಾಡು}} (\emph{alugāḍu}) is native Dravidian: DEDR 240 lists Kannada \emph{alaku}, \emph{alugu}, \textquotedblleft{}to move about, to shake, to tremble\textquotedblright{}, here with \kn{ಆಡು}.
\end{cousinweb}

\begin{grammarlens}[title={What you've built}]
One more word for getting around.
\end{grammarlens}

\subsection*{Your turn}
\begin{itemize}
\raggedright
  \item \emph{Say it:} \emph{alugāḍu}
  \item \emph{Say it:} \emph{alugāḍu}, once more
  \item \emph{Say it:} \emph{haḷeya bas kallina rasteyalli alugāḍuttittu.} — The old bus was rattling along the stony road.
  \item \emph{Recall:} say the Kannada for to park a vehicle, then the Kannada for to turn round, then say \emph{alugāḍu} again
\end{itemize}

\begin{tcolorbox}[breakable,colback=teal!4,colframe=teal!35!black,title={Before you move on}]
What does \kn{ಅಲುಗಾಡು} mean? (\textquotedblleft{}To wobble; to shake\textquotedblright{}.) Say it once more.
\end{tcolorbox}
\section[\texorpdfstring{\kn{ಚಲಿಸು} (calisu)}{ಚಲಿಸು (calisu)}]{\kn{ಚಲಿಸು} (calisu) — to move; to be in motion}

\label{lesson:KA-C445-calisu}



\begin{quote}
Before the new one: say the Kannada for to turn round, then the Kannada for to wobble.
\end{quote}

\subsection*{You'll want to know: \kn{ಚಲಿಸು}}
\textbf{\kn{ಚಲಿಸು}} — \emph{calisu} — \textquotedblleft{}to move; to be in motion\textquotedblright{}.

\begin{quote}
\kn{ರೈಲು} \kn{ನಿಧಾನವಾಗಿ} \kn{ಚಲಿಸಲು} \kn{ಶುರು} \kn{ಮಾಡಿತು}.
\end{quote}

\begin{quote}
\emph{railu nidhānavāgi calisalu śuru māḍitu.}
\end{quote}

\begin{quote}
The train slowly began to move.
\end{quote}

Past (it …): \kn{ಚಲಿಸಿತು}, \emph{calisitu}.

\kn{ಚಲಿಸುವ} \kn{ಮೆಟ್ಟಿಲು}, moving stairs, is an escalator.

\begin{cousinweb}
\textbf{\kn{ಚಲಿಸು}} (\emph{calisu}) is built from the Sanskrit root \emph{cal}, \textquotedblleft{}to move, to stir\textquotedblright{} (CDIAL 4715, \emph{calati}), with the Kannada verb ending -\emph{isu}.
\end{cousinweb}

\begin{grammarlens}[title={What you've built}]
One more word for getting around.
\end{grammarlens}

\subsection*{Your turn}
\begin{itemize}
\raggedright
  \item \emph{Say it:} \emph{calisu}
  \item \emph{Say it:} \emph{calisu}, once more
  \item \emph{Say it:} \emph{railu nidhānavāgi calisalu śuru māḍitu.} — The train slowly began to move.
  \item \emph{Recall:} say the Kannada for to turn round, then the Kannada for to wobble, then say \emph{calisu} again
\end{itemize}

\begin{tcolorbox}[breakable,colback=teal!4,colframe=teal!35!black,title={Before you move on}]
What does \kn{ಚಲಿಸು} mean? (\textquotedblleft{}To move; to be in motion\textquotedblright{}.) Say it once more.
\end{tcolorbox}
\section[\texorpdfstring{\kn{ಸಮೀಪಿಸು} (samīpisu)}{ಸಮೀಪಿಸು (samīpisu)}]{\kn{ಸಮೀಪಿಸು} (samīpisu) — to approach; to draw near}

\label{lesson:KA-C445-samipisu}



\begin{quote}
Before the new one: say the Kannada for to wobble, then the Kannada for to move.
\end{quote}

\subsection*{You'll want to know: \kn{ಸಮೀಪಿಸು}}
\textbf{\kn{ಸಮೀಪಿಸು}} — \emph{samīpisu} — \textquotedblleft{}to approach; to draw near\textquotedblright{}.

\begin{quote}
\kn{ವಿಮಾನ} \kn{ನಗರವನ್ನು} \kn{ಸಮೀಪಿಸಿತು}.
\end{quote}

\begin{quote}
\emph{vimāna nagaravannu samīpisitu.}
\end{quote}

\begin{quote}
The plane drew near the city.
\end{quote}

Past (I …): \kn{ಸಮೀಪಿಸಿದೆ}, \emph{samīpiside}.

\kn{ಸಮೀಪ} means near: \kn{ಸಮೀಪದಲ್ಲಿ}, nearby.

\begin{cousinweb}
\textbf{\kn{ಸಮೀಪಿಸು}} (\emph{samīpisu}) is built from Sanskrit \emph{samīpa}, \textquotedblleft{}near\textquotedblright{} (CDIAL 13226), with the verb ending -\emph{isu}: to come near.
\end{cousinweb}

\begin{grammarlens}[title={What you've built}]
One more word for getting around.
\end{grammarlens}

\subsection*{Your turn}
\begin{itemize}
\raggedright
  \item \emph{Say it:} \emph{samīpisu}
  \item \emph{Say it:} \emph{samīpisu}, once more
  \item \emph{Say it:} \emph{vimāna nagaravannu samīpisitu.} — The plane drew near the city.
  \item \emph{Recall:} say the Kannada for to wobble, then the Kannada for to move, then say \emph{samīpisu} again
\end{itemize}

\begin{tcolorbox}[breakable,colback=teal!4,colframe=teal!35!black,title={Before you move on}]
What does \kn{ಸಮೀಪಿಸು} mean? (\textquotedblleft{}To approach; to draw near\textquotedblright{}.) Say it once more.
\end{tcolorbox}
